\documentclass[frenchb]{article}
\usepackage[T1]{fontenc}
\usepackage[latin1]{inputenc}
%Pour utilisation sous unix
%\usepackage[utf8]{inputenc}
%\usepackage[utf8x]{inputenc}
\usepackage{a4wide}
\usepackage{graphicx}
\usepackage{amssymb}
\usepackage{color}
\usepackage{babel}

\begin{document}

\begin{figure}[t]
\centering
\includegraphics[width=5cm]{inp_n7.png}
\end{figure}

\title{\vspace{4cm} \textbf{Manuel utilisateur - it�ration 1}}
\author{Issam Alouane\\Remi Dabrowski\\Adam Bouam\\Jean-Baptiste Prevost }
\date{\vspace{7cm} D�partement Sciences du Num�rique - Premi�re ann�e \\
2021-2022 }

\maketitle

\newpage
\tableofcontents
\listoffigures

\newpage
\section{Menu principal}
Une fois que l'application est lanc�e gr�ce � la classe MenuPrincipal.java il y a la possibilit� soit de quitter, soit d'acc�der aux diff�rentes fonctionnalit�s. A ce stade il est uniquement possible de se rendre dans le sous menu personnages.

\section{Sous menu personnage}
Une fois dans le menu personnage vous avez la possibilit� de cr�er un personnage avec le bouton Nouveau. Une interface devrait s'ouvrir en demandant un nom et une description et aussi en demandant de choisir une race parmis les cat�gories propos�es, afin d'initialiser le personnage avant d'avoir la possibilit� de modifier ses statistiques.\\

A ce stade, le joueur peut soit annuler la cr�ation en cliquant sur "Annuler", soit avancer vers l'�tape suivante en cliquant sur "Suivant". L�, il peut modifier les statistiques du personnage (initialis�es par d�faut � partie de sa race). Ainsi, le joueur peut annuler toute la proc�dure de cr�ation en cliquant sur "Annuler" ou il peut enregistrer son personnage en cliquant sur "Enregistrer" \\

Une fois le personnage cr�� et enregistr�, un cadre avec le personnage cr�� devrait apparaitre. En cliquant sur celui-ci vous aurez la possibilit� de supprimer ou de modifier ce personnage (Cette fonctionnalit� n'est pas encore disponible).

\section{Menu �quipement personnage}
Ce menu n'est pas reli� � l'application principale. Il faut donc lancer ce menu gr�ce � la classe FenetreEquiper.java. Dans ce menu, vous avez cinq onglets, un onglet g�n�ral qui affiche les statistiques g�n�rales. Les quatre autres onglets concernent chaque �quipement avec la possibilit� de les changer gr�ce aux sliders et au bouton de validation. 

\end{document} 